Students: Alex Liddawi and Tarun Ramesh

Title: Analyzing the Feasibility of a Solar Sail Mission to Analyze the Composition of Asteroid Belts 

Abstract
The asteroid belt is a region of space located roughly between the orbits of Mars and Jupiter. It consists of millions of rocky, irregularly shaped bodies, which range from small comets to dwarf planets. Today, only a tiny amount of these places have been thoroughly mapped and surveyed. While Earth-based systems offer exceptional tracking capabilities to monitor their position, they cannot accurately detect the chemical composition of these bodies. To address this, many private organizations and nations have launched missions to survey the belt. However, due to fuel limitations, these missions are often only capable of analyzing one or two asteroids at a time. We propose a probe that can survey multiple objects utilizing a solar sail, a large, reflective membrane that utilizes the mass that photons from the Sun carry, creating a gentle yet continuous thrust to propel a spacecraft forward. To that end, it does not require an external source of fuel, making it useful for visiting multiple asteroids over a long duration of time.
Through conceptualizing and designing our own mission by analyzing alternative methods of propulsion, orbital periods, cost-effectiveness, and legal implications, this project determined the viability of a solar sail expedition to the asteroid belt. Our proposed design was accompanied with a physical scaled-down replica created with CAD software. 

Index Terms
Solar sail, Pallas (asteroid), spacecraft propulsion, asteroid belt survey, space probe design 
